
\documentclass{article}
\usepackage{colortbl}
\usepackage{makecell}
\usepackage{multirow}
\usepackage{supertabular}

\begin{document}

\newcounter{utterance}

\centering \large Interaction Transcript for game `bbh', experiment `web\_of\_lies', episode 206 with llama{-}3.1{-}8b{-}instruct.
\vspace{24pt}

{ \footnotesize  \setcounter{utterance}{1}
\setlength{\tabcolsep}{0pt}
\begin{supertabular}{c@{$\;$}|p{.15\linewidth}@{}p{.15\linewidth}p{.15\linewidth}p{.15\linewidth}p{.15\linewidth}p{.15\linewidth}}
    \# & \multicolumn{2}{c}{Player} && \multicolumn{2}{c}{Game Master} \\
    \hline

    \theutterance \stepcounter{utterance}  
    & & & \multicolumn{4}{p{0.6\linewidth}}{
        \cellcolor[rgb]{0.9,0.9,0.9}{
            \makecell[{{p{\linewidth}}}]{
                \texttt{\tiny{[P1$\langle$GM]}}
                \texttt{Evaluate a random boolean function expressed as a word problem.} \\
\\ 
\texttt{Q: Question: Fidel tells the truth. Jerry says Fidel tells the truth. Vina says Jerry tells the truth. Millicent says Vina lies. Raymond says Millicent lies. Does Raymond tell the truth?} \\
\texttt{A: Yes} \\
\\ 
\texttt{Q: Question: Kristian lies. Millie says Kristian lies. Maybelle says Millie tells the truth. Fidel says Maybelle lies. Leda says Fidel lies. Does Leda tell the truth?} \\
\texttt{A: Yes} \\
\\ 
\texttt{Q: Question: Kristian tells the truth. Michaela says Kristian lies. Raymond says Michaela tells the truth. Osvaldo says Raymond tells the truth. Jamey says Osvaldo tells the truth. Does Jamey tell the truth?} \\
\texttt{A: No} \\
\\ 
\texttt{Q: Question: Vernell lies. Conception says Vernell lies. Kandi says Conception lies. Michael says Kandi tells the truth. Osvaldo says Michael tells the truth. Does Osvaldo tell the truth?} \\
\texttt{A:} \\
            }
        }
    }
     \\ \\

    \theutterance \stepcounter{utterance}  
    & \multicolumn{4}{p{0.6\linewidth}}{
        \cellcolor[rgb]{0.561,0.737,0.561}{
            \makecell[{{p{\linewidth}}}]{
                \texttt{\tiny{[P1$\rangle$GM]}}
                \texttt{To evaluate the given boolean function, we can represent it as a truth table. However, since the problem is expressed as a word problem, we can analyze it step by step.} \\
\\ 
\texttt{Given: Vernell lies. Conception says Vernell lies. Kandi says Conception lies. Michael says Kandi tells the truth. Osvaldo says Michael tells the truth.} \\
\\ 
\texttt{1. Vernell lies, so Conception's statement is true (since Conception says Vernell lies).} \\
\texttt{2. Conception's statement is true, so Kandi's statement is false (since Kandi says Conception lies).} \\
\texttt{3. Kandi's statement is false, so Michael's statement is true (since Michael says Kandi tells the truth).} \\
\texttt{4. Michael's statement is true, so Osvaldo's statement is true (since Osvaldo says Michael tells the truth).} \\
\\ 
\texttt{Since Osvaldo's statement is true and Osvaldo says Michael tells the truth, and Michael says Kandi tells the truth, and Kandi says Conception lies, and Conception says Vernell lies, and Vernell lies, we can conclude that Osvaldo tells the truth.} \\
\\ 
\texttt{Therefore, the answer is: Yes} \\
            }
        }
    }
    & & \\ \\

    \theutterance \stepcounter{utterance}  
    & & & \multicolumn{2}{p{0.3\linewidth}}{
        \cellcolor[rgb]{0.9,0.9,0.9}{
            \makecell[{{p{\linewidth}}}]{
                \texttt{\tiny{[GM$|$GM]}}
                \texttt{Target: No} \\
            }
        }
    }
    & & \\ \\

    \theutterance \stepcounter{utterance}  
    & & & \multicolumn{2}{p{0.3\linewidth}}{
        \cellcolor[rgb]{0.9,0.9,0.9}{
            \makecell[{{p{\linewidth}}}]{
                \texttt{\tiny{[GM$|$GM]}}
                \texttt{game\_result = LOSE} \\
            }
        }
    }
    & & \\ \\

\end{supertabular}
}

\end{document}
